% !TeX program = xelatex
% 北京邮电大学《神经网络与深度学习（理论）》A2 课程实验报告
% 模板沿用 A1 报告工程（BUPTReport.sty），来源见 TEMPLATE_CREDITS.md。
\documentclass[12pt,a4paper,UTF8,fontset=fandol]{ctexart}
% 封面字段：仅需修改本文件，无需改动样式代码。
\newcommand{\ReportUniversity}{北京邮电大学}
\newcommand{\ReportCollege}{人工智能学院}
\newcommand{\ReportCourse}{神经网络与深度学习（理论）}
\newcommand{\ReportHeader}{神经网络与深度学习}
\newcommand{\ReportType}{课程实验报告}
\newcommand{\ReportAssignment}{A2 卷积神经网络图像分类}
\newcommand{\ReportTitle}{基于卷积神经网络的 CIFAR-10 图像分类实验}
\newcommand{\ReportAuthor}{罗开炜}
\newcommand{\ReportStudentID}{2024211735}
\newcommand{\ReportClass}{2024219112}
\newcommand{\ReportDate}{2026 年 11 月}

\usepackage{BUPTReport}
\makeatletter
\newcommand{\ReportBibControl}{\immediate\write\@auxout{\string\citation{report-bib}}}
\makeatother
\ExplSyntaxOn
\NewDocumentCommand{\ReportCompactURL}{m}
  {\group_begin:
   \tl_set:Nn \l_tmpa_tl {#1}
   \regex_replace_once:nnN {\A https?://(www\.)?} {} \l_tmpa_tl
   \href{#1}{\mbox{\exp_args:NV \nolinkurl \l_tmpa_tl}}
   \group_end:}
\ExplSyntaxOff

\hypersetup{pdfsubject={北京邮电大学 A2 卷积神经网络图像分类课程实验报告}}
\hypersetup{pageanchor=false}

\begin{document}
\makebuptcover
\hypersetup{pageanchor=true}

\frontmatterbupt
\begin{center}
    {\heiti\zihao{3} 摘\quad 要}
\end{center}
\vspace{0.8em}
构建CIFAR-10卷积分类器，通过107次正式训练比较Dropout、归一化及超参数选型，并验证深度退化、校准和平移敏感性的改进。基准三种子验证准确率为$83.15\pm0.29\%$，FC Dropout0.7和BN分别为$85.76\pm0.16\%$与$86.77\pm0.16\%$。随机搜索与五折选型后的全量重训测试89.25\%；增强、平滑和加宽单模型测试93.87\%，三模型集成94.36\%。残差连接、温度缩放及抗混叠池化分别改善相应问题。

\vspace{1.3em}
\noindent\textbf{关键词：}\hspace{0.5em}卷积神经网络；CIFAR-10；图像分类；Dropout；归一化
\clearpage

\tableofcontents
\clearpage

\mainmatterbupt
\section{实验目的与任务}
CIFAR-10包含飞机、汽车、鸟等10类低分辨率自然图像\cite{krizhevskyCifar}。本实验围绕课程A2要求，构建可训练的卷积神经网络，研究Dropout与多种归一化对泛化的影响，通过随机搜索和分层交叉验证选择超参数，并对性能提升与网络局限开展验证。

卷积通过局部连接与参数共享利用图像的空间结构\cite{lecun1998gradient}。与将图像展平的多层感知机相比，它提供了更合适的结构先验。本报告据此区分三个问题：基础CNN能否稳定学习；正则化与归一化在固定预算下如何改变验证表现；在提高分类准确率之外，能否改善深度退化、校准和平移稳定性。

实验共定义107次正式训练，包含多随机种子对照、24组随机搜索、四候选的五折验证及最终重训。作者此前A1报告的NumPy MLP测试准确率50.39\%作为历史参考。

\section{模型原理与方法}
\subsection{卷积、池化与分类目标}
设输入特征为$x$，卷积层第$k$个输出通道为
\begin{equation}
 z_{k,i,j}=b_k+\sum_c\sum_{u,v}W_{k,c,u,v}x_{c,i+u,j+v},\qquad a=\max(0,z).
\end{equation}
小卷积核逐层组合局部特征，池化降低空间尺寸。本文使用3$\times$3卷积和ReLU，初始化采用针对整流网络的Kaiming方法\cite{he2015rectifiers}。分类器产生10维logits $z$，softmax输出$p_k=\exp(z_k)/\sum_j\exp(z_j)$，普通训练目标为交叉熵$-\log p_y$。步幅下采样和边界处理会影响分类器的平移稳定性。

\subsection{Dropout与归一化}
Dropout在训练时以概率$p$丢弃激活，并对保留值缩放：
\begin{equation}
 \widetilde h=\frac{m\odot h}{1-p},\qquad m_i\sim\operatorname{Bernoulli}(1-p).
\end{equation}
评估时关闭随机丢弃。该方法限制特征间的共同适应，可用于缓解过拟合\cite{srivastava2014dropout}。本文分别在全连接隐藏层之后使用元素Dropout，在卷积段池化之后使用整通道Dropout2d。

各归一化方法都可概括为
\begin{equation}
 \widehat x_i=\frac{x_i-\mu_{\mathcal S}}{\sqrt{\sigma^2_{\mathcal S}+\epsilon}},\qquad
 y_i=\gamma_c\widehat x_i+\beta_c,
\end{equation}
差别主要在统计集合$\mathcal S$及评估时统计量的处理。对NCHW张量，本文采用的具体变体见表\ref{tab:normdef}。
\begin{table}[htbp]\centering\small
\caption{本文归一化实现的统计维度与仿射参数}\label{tab:normdef}
\begin{tabularx}{\linewidth}{llX}
\toprule 方法 & 统计范围 & 实现与评估方式\\\midrule
BN & 同一通道的$N,H,W$ & BatchNorm2d；评估使用运行均值与方差\\
LN & 单一样本的$C,H,W$ & GroupNorm(1,C)，仿射参数仅逐通道\\
IN & 单一样本单一通道的$H,W$ & GroupNorm(C,C)，不使用运行统计\\
GN & 单一样本每组通道及$H,W$ & GroupNorm(8,C)，不依赖批统计\\
\bottomrule\end{tabularx}\end{table}
\FloatBarrier
BN、LN、IN与GN分别依据相应文献\cite{ioffe2015batchnorm,ba2016layernorm,ulyanov2016instancenorm,wu2018groupnorm}。BN与Dropout的相互作用与插入位置及统计偏移有关\cite{li2019disharmony}。

\subsection{性能提升及问题改进}
随机裁剪和翻转扩展训练时的图像变化；Cutout随机遮挡局部区域\cite{devries2017cutout}。标签平滑以$q_k=(1-\varepsilon)\mathbf1[k=y]+\varepsilon/10$替换硬标签，本文取$\varepsilon=0.1$\cite{szegedy2016rethinking}。增宽网络增加特征容量，残差块则用$h_{l+1}=\operatorname{ReLU}(F(h_l)+S(h_l))$改善深层优化；$S$在尺寸变化时为投影捷径\cite{he2016resnet}。

分类准确率只依赖最大概率类别，校准则关注概率与实际正确率是否匹配。负对数似然（NLL）与15等宽置信度分箱的期望校准误差（ECE）分别为
\begin{equation}
 \operatorname{NLL}=-\frac1n\sum_i\log p_{i,y_i},\qquad
 \operatorname{ECE}=\sum_{b=1}^{15}\frac{|B_b|}{n}|\operatorname{acc}(B_b)-\operatorname{conf}(B_b)|.
\end{equation}
温度缩放用$\operatorname{softmax}(z/T)$调整概率，$T>0$由验证NLL选择\cite{guo2017calibration}。正温度不改变argmax，因而不改变分类准确率。抗混叠下采样则在降采样前加入低通滤波\cite{zhang2019shift}，本文用它检验CNN对小幅输入平移的稳定性。

\section{实现与训练流程}
\subsection{主要模块与数据管线}
\texttt{config.py}与\texttt{experiments.py}定义训练配置和实验分组；\texttt{data.py}处理数据读取、划分与预处理；\texttt{models.py}实现CNN、Plain/ResNet及池化模块；\texttt{train.py}完成训练、评估与保存；\texttt{analysis.py}和\texttt{report.py}负责指标汇总与图表生成。

数据从CIFAR-10的Python批文件读取。DeviceLoader将图像和标签放到目标设备，按批完成采样、标准化和训练增强。每轮依次执行前向传播、交叉熵计算、反向传播、SGD更新及学习率调度，随后切换为评估模式计算指标。

\subsection{评估、保存与运行环境}
报告中的训练准确率来自固定5000张训练子集，在评估模式下计算；泛化差距为其与验证准确率之差，以百分点（pp）表示。模型选择和最终评估统一使用最后一轮权重。每轮保存训练断点，任务结束后保存配置、指标历史、模型权重及预测logits。

实验使用单张RTX 5090（约32GB显存），Ubuntu 22.04、Python 3.12.3、PyTorch 2.8.0+cu128和NumPy 2.3.2；训练采用BF16混合精度前向与FP32损失。基础对照固定随机种子0、1、2，其余实验设置见表\ref{tab:design}。

\section{实验设置}
\subsection{数据划分与预处理}
CIFAR-10有50000张训练图和10000张测试图，每张为32$\times$32的RGB图像，共10类\cite{krizhevskyCifar}。基础实验沿用A1分层划分：45000张训练、5000张验证，种子20261011，每类分别4500与500张。五折划分种子20261102，每折每类1000张，单折训练40000张、验证10000张。最终选定配置另在全部50000张训练图上重训。

像素先除255，再按通道标准化。通道均值与标准差分别为
\[ \boldsymbol\mu=(0.4914,0.4822,0.4465),\qquad \boldsymbol\sigma=(0.2470,0.2435,0.2616). \]
标准化常数取自完整训练池。增强时补4像素边界，随机裁剪回32$\times$32并以0.5概率水平翻转；填充值为标准化后的零。Cutout边长为16，中心随机。

\subsection{基准网络与优化参数}
基准结构见图\ref{fig:architecture}。三个卷积段各含两层3$\times$3卷积，通道为64/128/256；三次池化将空间尺寸降至4$\times$4，分类器为4096$\to$256$\to$10。无归一化模型有2196810个参数；归一化版本移除卷积bias并加入affine，共2197706个参数。
\begin{figure}[htbp]\centering
\includegraphics[width=\linewidth]{architecture.pdf}
\caption{基准CNN的张量变化与模块位置。叠片示意通道组，尺寸为通道$\times$高$\times$宽。}\label{fig:architecture}
\end{figure}

默认训练30轮，batch=128，峰值学习率0.02，采用SGD、动量0.9、Nesterov及weight decay=$5\times10^{-4}$。前2轮线性预热，之后余弦退火至0。基础实验使用原图和硬标签，增强与标签平滑在B1中逐项加入。

\subsection{实验矩阵与超参数选择}
\begin{table}[htbp]\centering\small
\caption{正式实验设计。次数为独立训练任务数；“三种子”指0、1、2，其余为种子0。}\label{tab:design}
\begin{tabularx}{\linewidth}{lXrr}
\toprule 组 & 内容 & 次数 & 轮数\\\midrule
E0 & 无归一化、无Dropout基准，三种子 & 3 & 30\\
E1 & FC概率0.3/0.5/0.7；卷积0.1/0.2与FC0.5组合，三种子 & 15 & 30\\
E2归一化 & BN/LN/IN/GN，三种子 & 12 & 30\\
E2批大小 & BN/GN/LN，batch=16且学习率降至0.0025 & 3 & 30\\
E2学习率 & none/BN/GN与0.005/0.1/0.4组合 & 9 & 30\\
E2组合 & BN+FC0.5，BN+卷积0.1+FC0.5，三种子 & 6 & 30\\
E3搜索 & 24组随机超参数 & 24 & 15\\
E3交叉验证 & 搜索前4名，各做5折 & 20 & 30\\
E3最终 & CV均值最优配置，全50000张重训 & 1 & 30\\
B1 & 增强、Cutout、标签平滑、加宽、架构替换 & 6 & 40\\
B2深度 & Plain/ResNet各20/32/56层 & 6 & 30\\
B2平移 & MaxBlurPool，无增强/裁剪翻转 & 2 & 40\\\bottomrule
\end{tabularx}\end{table}

随机搜索种子20261103，从$[0.005,0.2]$和$[10^{-5},2\times10^{-3}]$按对数均匀抽取学习率与weight decay，从$\{0,0.3,0.5\}$、$\{64,128,256\}$和$\{\mathrm{BN},\mathrm{GN}\}$抽取FC Dropout、batch和归一化。24组各训练15轮，按末轮验证准确率选前4名；随后从头训练每个候选的五折模型，以折均值选配置，最后在全部训练池重训30轮。

B1在CV最优配置上做40轮累计配方比较，依次加入裁剪翻转、Cutout、标签平滑和VGG加宽（128/256/512通道），最后一项为ResNet-56架构替换，采用16/32/64通道及全局平均池化分类头。B2深度实验统一BN、学习率0.1与裁剪翻转，比较同主路径深度有无捷径。两类集成均为softmax概率平均：一类由CV赢家的五个折模型组成，另一类从B1按验证准确率选3个成员。

\subsection{评价指标与模型选择}\label{sec:boundaries}
Accuracy衡量分类正确率，NLL衡量真实类别的预测概率，ECE衡量置信度与正确率的偏差。Accuracy与ECE用百分比表示，NLL取自然对数。三种子实验报告均值$\pm$样本标准差（ddof=1），五折验证报告折均值与折间标准差，其余为单种子结果。

随机搜索按验证准确率选候选，五折均值用于内部选型；集成成员和温度参数也由验证数据确定。全部正式训练完成后，统一进行完整测试集评估。

\Needspace{12\baselineskip}
\section{实验结果与分析}
\subsection{基准CNN与Dropout}
\textbf{无增强基准存在明显过拟合。}三种子末轮验证准确率为$83.15\pm0.29\%$，固定训练子集的评估准确率为$100.00\pm0.00\%$，泛化差距$16.85\pm0.29$个百分点。图\ref{fig:baseline}中，训练曲线逐渐饱和，验证曲线仍明显低于训练曲线。
\begin{figure}[htbp]\centering
\includegraphics[width=\linewidth]{baseline_curves.pdf}
\caption{基准、FC Dropout 0.5与BN的学习曲线（种子0）。实线为固定5000张训练子集的评估准确率，虚线为验证准确率。}\label{fig:baseline}
\end{figure}
\FloatBarrier
\begin{table}[htbp]\centering\small
\caption{Dropout对照：三种子均值$\pm$标准差，验证损失为三种子均值。}\label{tab:dropout}
\ReportTable{tables/dropout}
\end{table}

表\ref{tab:dropout}中，FC Dropout由0增至0.7时，验证准确率逐步提高；0.7配置比基准高2.61个百分点。随机丢弃限制了隐藏特征间的共同适应，缩小了训练与验证差距。卷积0.1+FC0.5为$85.23\pm0.18\%$，与单用FC0.5的表现接近。卷积0.2组合的验证损失降至0.598，准确率为85.09\%，说明预测概率的改善与类别判断的改善并不同步。
\begin{figure}[htbp]\centering
\includegraphics[width=\linewidth]{dropout_gap.pdf}
\caption{Dropout的验证表现。浅色点为三个种子，菱形及误差条为均值$\pm$标准差。}\label{fig:dropout}
\end{figure}

\FloatBarrier
\subsection{归一化及其与训练配方的相互作用}
表\ref{tab:norm}在固定学习率0.02、batch128、30轮条件下比较归一化。BN验证均值86.77\%，比无归一化高3.62个百分点；GN、LN、IN分别为85.95\%、85.35\%、84.52\%。归一化使中间激活的尺度更稳定，改善了当前配方的收敛与泛化；其中BN的验证均值最高。
\begin{table}[htbp]\centering\small
\caption{归一化对照：三种子均值$\pm$标准差。}\label{tab:norm}
\ReportTable{tables/norm_report}
\end{table}
\begin{figure}[htbp]\centering
\includegraphics[width=\linewidth]{norm_curves.pdf}
\caption{各归一化方法的验证收敛曲线（种子0）。}\label{fig:norm}
\end{figure}

将batch从128改为16，并将学习率线性降至0.0025后，三种归一化的验证变化均在0.3个百分点以内（表\ref{tab:batch}）。BN在这组小批量配方下仍保持相近表现；GN与LN的统计量按样本计算，不依赖批内样本数。
\begin{table}[htbp]\centering\small
\caption{批大小配方比较，均为种子0；batch16同时采用八分之一学习率。}\label{tab:batch}
\begin{tabular}{lccc}
\toprule
归一化 & Batch 128 (\%) & Batch 16 (\%) & 变化 (pp) \\
\midrule
BatchNorm & 86.68 & 86.86 & +0.18 \\
GroupNorm & 85.76 & 86.02 & +0.26 \\
LayerNorm & 85.82 & 85.66 & -0.16 \\
\bottomrule
\end{tabular}

\end{table}

学习率对照见表\ref{tab:lr}及图\ref{fig:lr}。峰值0.1时none/BN/GN分别为85.22\%、89.00\%、87.28\%；峰值0.4配方下，无归一化网络在第1轮预热阶段出现NaN，而BN与GN仍得到88.50\%与84.82\%。归一化拓宽了本配方的可训练学习率范围；GN在0.4下准确率回落，表明步长增大后优化效果也会下降。
\begin{table}[htbp]\centering\small
\caption{不同峰值学习率的末轮验证准确率（\%），种子0；0.02列复用基础对照。}\label{tab:lr}
\begin{tabular}{lcccc}
\toprule
归一化 & lr=0.005 & lr=0.02 & lr=0.1 & lr=0.4 \\
\midrule
无 & 79.16 & 83.46 & 85.22 & 发散 \\
BatchNorm & 81.98 & 86.68 & 89.00 & 88.50 \\
GroupNorm & 81.72 & 85.76 & 87.28 & 84.82 \\
\bottomrule
\end{tabular}

\end{table}
\begin{figure}[htbp]\centering
\includegraphics[width=\linewidth]{norm_lr.pdf}
\caption{学习率敏感性（种子0）。无归一化的0.4任务在第1轮出现NaN。}\label{fig:lr}
\end{figure}

表\ref{tab:combo}中，BN加FC Dropout0.5由$86.77\pm0.16\%$提高至$88.72\pm0.45\%$，再加卷积Dropout0.1则为$87.83\pm0.08\%$。BN与FC Dropout在本组实验中形成互补，而进一步增加卷积通道丢弃降低了验证均值。卷积层和分类器上的正则化强度需要匹配。
\begin{table}[htbp]\centering\small
\caption{Dropout与BN组合的末轮验证准确率（\%），三种子均值$\pm$样本标准差。}\label{tab:combo}
\begin{tabular}{lccc}
\toprule
归一化 & 无 Dropout & FC 0.5 & Conv 0.1 + FC 0.5 \\
\midrule
无 & 83.15 $\pm$ 0.29 & 85.07 $\pm$ 0.19 & 85.23 $\pm$ 0.18 \\
BatchNorm & 86.77 $\pm$ 0.16 & 88.72 $\pm$ 0.45 & 87.83 $\pm$ 0.08 \\
\bottomrule
\end{tabular}

\end{table}

\FloatBarrier
\Needspace{7\baselineskip}
\subsection{搜索、五折选型与最终测试}
24组搜索结果完整列于附录表\ref{tab:search}。前4名为cfg00、cfg02、cfg08、cfg18，15轮验证准确率分别为89.42\%、88.64\%、88.42\%、88.34\%。cfg01采用GN、学习率0.1578和FC Dropout0.5，损失有限但末轮验证仅10.00\%，表现为优化塌缩。
\begin{table}[htbp]\centering\footnotesize
\caption{四候选的五折验证准确率（\%）：30轮、种子0，报告折间标准差。}\label{tab:cv}
\begin{tabular}{lcccccc}
\toprule
配置 & 折1 & 折2 & 折3 & 折4 & 折5 & 均值 $\pm$ STD \\
\midrule
cfg00 & 88.73 & 88.53 & 88.50 & 88.22 & 88.57 & 88.51 $\pm$ 0.18 \\
cfg08 & 88.42 & 88.39 & 88.39 & 88.06 & 88.68 & 88.39 $\pm$ 0.22 \\
cfg02 & 87.95 & 87.78 & 88.55 & 88.04 & 88.33 & 88.13 $\pm$ 0.31 \\
cfg18 & 87.58 & 87.53 & 87.56 & 87.42 & 87.76 & 87.57 $\pm$ 0.12 \\
\bottomrule
\end{tabular}

\end{table}

表\ref{tab:cv}中cfg00均值$88.51\pm0.18\%$，按预定规则被选中；cfg08为$88.39\pm0.22\%$，均值差为0.12个百分点。cfg00的精确配置为BN、batch128、学习率0.12089、weight decay=0.0009851、FC Dropout0.3。
\begin{table}[htbp]\centering\small
\caption{最终测试结果。折模型行报告五模型均值$\pm$标准差；A1为历史参考。}\label{tab:final}
\ReportTable{tables/final}
\end{table}

在全部50000张训练图上重训后，单模型测试89.25\%、NLL0.405；五折模型概率集成测试90.20\%、NLL0.319，较各折单模型的均值88.28\%更高。概率平均综合了不同折模型的判断，降低单个模型预测误差的影响，同时增加推理成本。
\FloatBarrier

\Needspace{16\baselineskip}
\section{附加实验：性能提升与问题改进}
\subsection{增强、正则化、加宽与集成}
B1以CV最优配置为起点，固定45000/5000划分、种子0和40轮预算，逐项加入增强、平滑与加宽，最后替换架构，结果见表\ref{tab:b1}与图\ref{fig:b1}。
\begin{table}[H]\centering\small
\caption{B1累计配方及架构替换，种子0、40轮。集成成员按验证表现选择。}\label{tab:b1}
\ReportTable{tables/b1_ablation}
\end{table}
\begin{figure}[H]\centering
\includegraphics[width=\linewidth]{b1_ablation.pdf}
\caption{B1相邻配方的验证准确率变化（种子0）。虚线分隔架构替换；绝对准确率与参数量见表\ref{tab:b1}。}\label{fig:b1}
\end{figure}

裁剪翻转将测试准确率从88.44\%提高至92.64\%，收益最大；它扩展了训练图像的空间变化。Cutout和标签平滑进一步改变正则化方式，准确率小幅波动。加宽VGG后，参数量由约2.20M增至6.68M，测试准确率达到93.87\%。ResNet-56约0.86M参数，继承VGG选出的优化配方，测试为91.63\%，低于加宽VGG。

按末轮验证准确率选出的三个成员为加宽VGG、标签平滑VGG和裁剪翻转VGG。三模型概率平均达到94.36\%测试准确率，比加宽单模型高0.49个百分点，集成综合了不同配方的预测。

\FloatBarrier
\Needspace{7\baselineskip}
\subsection{深度退化与残差连接}
对Plain和ResNet使用相同主路径深度、BN、裁剪翻转、峰值学习率0.1，训练30轮。图\ref{fig:depth}和表\ref{tab:depth}显示，Plain从20层增至56层时，训练子集误差由5.38\%升至17.20\%，验证误差由10.76\%升至19.58\%。训练与验证误差同时上升，呈现深度增加后的优化退化。
\begin{figure}[H]\centering
\includegraphics[width=\linewidth]{depth_train_error.pdf}
\caption{Plain与ResNet的训练误差曲线（种子0）。}\label{fig:depth}
\end{figure}
\begin{table}[H]\centering\small
\caption{不同深度的末轮错误率（\%），种子0。}\label{tab:depth}
\begin{tabular}{lcccc}
\toprule
层数 & Plain 训练误差 & Plain 验证误差 & ResNet 训练误差 & ResNet 验证误差 \\
\midrule
20 & 5.38 & 10.76 & 3.62 & 8.78 \\
32 & 7.68 & 11.76 & 2.80 & 7.76 \\
56 & 17.20 & 19.58 & 2.46 & 7.62 \\
\bottomrule
\end{tabular}

\end{table}

ResNet的相应训练误差由3.62\%降至2.46\%，验证误差由8.78\%降至7.62\%，支持残差捷径缓解深度退化\cite{he2016resnet}。残差捷径提供较直接的梯度传播路径，使深层网络更容易优化。

\FloatBarrier
\Needspace{8\baselineskip}
\subsection{置信度失真与温度缩放}
表\ref{tab:calibration}同时保留正确率、NLL和ECE。无增强基础模型平均置信度94.65\%，而准确率88.44\%，显示过度自信。Cutout原始ECE为1.98\%，验证集拟合$T=1.26$后测试ECE降至0.54\%，NLL从0.214降至0.205，而准确率仍为93.01\%。图\ref{fig:reliability}提供分箱观察。
\begin{table}[H]\centering\small
\renewcommand{\arraystretch}{1.0}
\caption{测试集校准结果（种子0）。TS为温度缩放，Aug.为裁剪翻转，LS为标签平滑；Acc./Conf./ECE单位为\%。}\label{tab:calibration}
\begin{tabular}{lccccccc}
\toprule
模型 & Acc. & Conf. & NLL & ECE & $T$ & NLL(TS) & ECE(TS) \\
\midrule
Base & 88.44 & 94.65 & 0.441 & 6.26 & 1.69 & 0.356 & 1.69 \\
Aug. & 92.64 & 96.09 & 0.251 & 3.47 & 1.55 & 0.215 & 0.84 \\
Cutout & 93.01 & 94.92 & 0.214 & 1.98 & 1.26 & 0.205 & 0.54 \\
LS & 92.84 & 84.97 & 0.281 & 7.87 & 0.63 & 0.222 & 1.52 \\
Wide & 93.87 & 85.90 & 0.254 & 8.01 & 0.63 & 0.192 & 1.25 \\
ResNet56 & 91.63 & 85.43 & 0.312 & 6.20 & 0.71 & 0.271 & 2.01 \\
\bottomrule
\end{tabular}

\end{table}
\begin{figure}[H]\centering
\includegraphics[width=.92\linewidth]{reliability.pdf}
\caption{可靠性图（15等宽箱）。上排为箱内平均置信度与准确率，下排为样本占比；虚线表示理想校准。}\label{fig:reliability}
\end{figure}

\textbf{标签平滑模型表现为欠自信。}LS模型置信度84.97\%低于准确率92.84\%（ECE为7.87\%）；加宽平滑模型也欠自信。平滑目标降低了正确类别的目标概率。两者拟合$T\approx0.63$后，logits差异被放大，ECE分别降至1.52\%、1.25\%。

\FloatBarrier
\Needspace{7\baselineskip}
\subsection{平移敏感性与抗混叠池化}
测试对$(d_y,d_x)\in[-4,4]^2$的81种整数位移进行评估，采用标准化零填充和边界截断。MaxBlurPool先做步幅1的2$\times$2最大池化，再用$[1,2,1]^\mathsf{T}[1,2,1]/16$的3$\times$3逐通道核低通滤波、反射padding并以步幅2下采样。

设图像$i$在81个位置中被预测为类别$c$的次数为$n_{i,c}$，一致性为不同位置对预测相同的比例，再对图像取平均：
\begin{equation}
 C=\frac1n\sum_{i=1}^n\frac{\sum_c n_{i,c}(n_{i,c}-1)}{81\times80}.
\end{equation}

\begin{table}[H]\centering\small
\caption{池化与增强的平移评估（种子0，40轮）。平均包含零位移，共81个位置。}\label{tab:shift}
\ReportTable{tables/shift}
\end{table}
\begin{figure}[H]\centering
\includegraphics[width=\linewidth]{shift_accuracy.pdf}
\caption{按$\max(|d_y|,|d_x|)$分组的平均准确率。半径0至4分别含1、8、16、24、32个位移。}\label{fig:shift}
\end{figure}

\textbf{抗混叠主要改善扰动稳定性。}无增强时平移平均准确率由83.06\%升至85.88\%，与裁剪翻转结合后也保持收益（表\ref{tab:shift}）。低通滤波减轻了下采样对小幅平移的敏感性。

\Needspace{30\baselineskip}
\section{结论}
本实验完成了CNN构建、Dropout与归一化对照及交叉验证选型，并通过增强、加宽和集成提升分类性能。进一步围绕深度退化、置信度失准和平移敏感开展改进实验，形成从模型训练、泛化分析到问题验证的完整实践。各项任务的完成内容与主要结论汇总于表\ref{tab:requirements}。

\begin{table}[H]\centering\small
\linespread{1.10}\selectfont
\renewcommand{\arraystretch}{1.12}
\caption{作业任务、完成内容与主要结论}\label{tab:requirements}
\begin{tabularx}{\linewidth}{@{}p{2.8cm}X p{1.8cm}@{}}
\toprule
作业任务 & 完成工作与主要结果 & 正文证据\\
\midrule
\textbf{必做：CNN分类} & 实现卷积、池化、分类层及训练流程；基准三种子验证准确率为$83.15\pm0.29\%$。 & 图\ref{fig:architecture}\newline 表\ref{tab:dropout}\\
\addlinespace[5pt]
\textbf{必做：Dropout} & 比较FC概率0.3/0.5/0.7及卷积Dropout组合；FC~0.7验证准确率为$85.76\pm0.16\%$，较基准提高2.61个百分点，缩小训练与验证差距。 & 表\ref{tab:dropout}\\
\addlinespace[5pt]
\textbf{必做：归一化} & 比较BN/LN/IN/GN及批大小、学习率和Dropout组合。BN验证准确率为$86.77\pm0.16\%$，BN与FC~0.5组合为$88.72\pm0.45\%$。 & 表\ref{tab:norm}--\ref{tab:combo}\\
\addlinespace[5pt]
\textbf{必做：交叉验证} & 24组随机搜索后，对前4名做五折选型；选定配置在全部50000张训练图上重训，测试准确率为89.25\%。 & 表\ref{tab:cv}、\ref{tab:final}\\
\midrule
\textbf{附加一：}\newline\textbf{性能提升} & 比较增强、Cutout、标签平滑、加宽及集成。裁剪翻转使同组测试准确率由88.44\%升至92.64\%；加宽模型与三模型集成分别为93.87\%、94.36\%。 & 表\ref{tab:b1}\\
\addlinespace[5pt]
\textbf{附加二：}\newline\textbf{问题分析与改进} & 深度退化：残差连接将56层网络验证错误率由19.58\%降至7.62\%。\newline
置信度失准：温度缩放将Cutout模型ECE由1.98\%降至0.54\%。\newline
平移敏感：无增强时，抗混叠池化将平移平均准确率由83.06\%提高至85.88\%。 & 表\ref{tab:depth}--\ref{tab:shift}\\
\bottomrule
\end{tabularx}
\end{table}

部分比较采用单种子，批大小对照同时调整了学习率，B1的架构替换也改变了模型容量。搜索和五折验证共用训练池，固定验证集重复参与方法比较与温度拟合。准备阶段冒烟曾评估测试集前500张；正式结果来自全部训练完成后的完整测试集分析。


\FloatBarrier
\clearpage
\phantomsection
\addcontentsline{toc}{section}{参考文献}
\bibliographystyle{gbt7714-numeric}
\ReportBibControl
\begingroup
\linespread{1.12}\small
\raggedright
\let\url\ReportCompactURL
\interlinepenalty=10000
\setlength{\bibsep}{2pt}
\bibliography{bibliography_style,references}
\endgroup

\clearpage
\appendix
\ctexset{section={name={附录\space,},number=\Alph{section}}}
\section{复现信息与补充结果}
\subsection{源码与复现}
项目仓库：\url{https://github.com/kaikailuo/cifar10_cnn}。按仓库README安装依赖并放置数据后，运行下列命令。实验配置与逐轮结果位于\path{results/}，汇总见\path{results/analysis/analysis.json}。
\begin{quote}\small\ttfamily
python -m src.run --jobs 4
\end{quote}

\subsection{作者贡献与AI辅助}
本项目使用AI工具辅助代码实现、实验运行管理、结果整理与报告撰写。实验目标、方案选择及关键技术决策由作者确定，作者负责审阅实验结果与最终提交内容。

\subsection{随机搜索完整结果}
\begin{center}\footnotesize\linespread{1}\selectfont
\renewcommand{\arraystretch}{1.05}
\captionsetup{type=table}
\caption{24组随机搜索末轮验证结果（种子0，15轮），按验证准确率排序。}\label{tab:search}
\ReportTable{tables/search}
\end{center}

\end{document}
